\chapter{Introduction}
\label{chap:intro}

The decentralized administration of the Internet is one of the basic enablers of the rapid growth of the Internet as a whole, thanks to the autonomous nature of different network systems~\cite{why-internet-works}.
Different Internet Service Providers (ISPs) play a critical role in the functionality of the Internet, as they connect their end customers (Internet users) to the Internet backhaul.
The power dynamic of this system inherently gives control over shaping consumption patterns to the ISPs: consolidation of power at the hands of a few gives the ISPs the unique ability to control, throttle, or monitor the flow of information, in addition to enabling oppressive practices such as censorship and surveillance of end users. 
Governmental control of the telecommunication infrastructure often leads to the repression of communities during times of political unrest~\cite{digital-repression-palestine,internet-shutdown-human-rights}, severely limiting the possibility for social actors to dissent and organize themselves.
In addition to policial repression engaged by the government, physical distruptions can also limit free speech.
Most recently, communication disruptions due to airstrikes on main fiber optic lines, and the depletion of fuel for generators, have severely limited the information flow to the population of Gaza since the start of the Israel-Hamas war in October 2023, halting humanitarian coordination, and leaving the affected population entirely cut off from life-saving information and assistance \cite{gaza-brief,gaza-report}.

Together, the aforementioned practices have led to a growing interest in peer-to-peer (p2p) messaging platforms, decentralized social networks~\cite{bluesky,mastodon}, and distributed ledger-based applications~\cite{bitcoin,ethereum}.
To guarantee free speech in disaster zones, during protests, and for citizens living under oppressive regimes, we decide to focus on unstructured p2p networks, as they inherently exhibit a higher level of resilience to topology changes and churn, in comparison to structured p2p networks.
Structured p2p networks usually require an expensive randomized approach to establish resilient topologies, as they need to spread information across random peers to prevent dependence on any single peer~\cite{guzman22, guzman25, kademlia}. Additionally, iterative diffusion of the information causes data plane overhead~\cite{guzman25}, which increases the message propagation latency~\cite{guzman22, guzman23}.
Popular blockchains, which use a structued p2p overlay, employ concentration of services in a few highly available nodes, poses a threat to resilience despite a high reliability of those nodes, as a single network partition can prevent the entire network from communicating~\cite{guzman24}.
% Other famous structured p2p networks include systems as Chord~\cite{chord}, Kademlia~\cite{kademlia}, and the InterPlanetary File System (IPFS)~\cite{ipfs}.
Unstructured p2p networks, on the other hand, are more suitable for highly-dynamic alternative communication methods, often necessary to guarantee free decentralized communication means, as they usually do not distribute central information amongst the participating nodes.
Most notably, the rise of unstructured p2p messaging platforms like Firechat~\cite{firechat}, Briar~\cite{briar}, bitchat~\cite{bitchat}, and Amigo~\cite{amigo}, have proven that there is a growing shift in user interest patterns away from traditional Centralized Social Networks towards diverse p2p social networks.

After exploring the different types of p2p overlay networks, we discuss the choice of the communication protocol, which dictates the limited physical device communication radius.
We chose to focus on a stable protocol that most modern devices support and the populace is most likely to use in disaster zones: Bluetooth, and more specifically Bluetooth Low Energy (BLE)~\cite{bt-spec}.
We decided to focus on Bluetooth Low Energy (BLE) as it is more power efficient, over classic Basic Rate/Enhanced Data Rate (BR/EDR), as BLE conserves device battery life, it allows users to communicate over short ranges over a longer period of time than comparable communication protocols e.g. WiFi Direct~\cite{wifi-ble-power}, %which is critical to save lives.

Next, we discuss the bottleneck of unstructured p2p networks: the communication algorithm.
Most unstructured p2p networks rely on gossip protocols to disseminate information across the network, which reduces its performance and poses a threat to resilience (cite guided research paper?).
Gossip protocols usually require constant chatter between peers to reach consensus about the system state, which introduces a massive performance overhead.
This performance overhead poses a risk to usablity of such systems in disaster zones, as peer devices often have limited connectivity and power supply, which together with the gossip protocol, can severely limit the peers' ability to communicate.
The introduction of permissioned peers, such as high-availability bootstrapping nodes or signaling servers, affect the network's ability to communicate in disaster zones, as those nodes are then often affected by network partitions, political censorship, or physical disruptions, which halts the entire network from communicating.
While such bootstrapping nodes ease the process of joining the network, discovering peers, and (potentially) disseminating state, they weaken the peers' ability to function without their presence, which is, as we show in this paper, oftentimes unnecessary.
Instead of using permissioned nodes to facilitate communication by achieving total order of transactions, we suggest that using a partial ordering of transactions using a distributed Directed Acyclic Graph Conflict Free Replicated Data Types (DAG-CRDT), in which each node contains blocks of data, as well as edges which reference previous blocks, lends itself well to storing persistent communication data needed for social networking.
This knowledge base, also known as the block frontier, provides proof to the local state of the node which created it, and iteratively ordering it on each participants' local device using a total order, achieves an eventual total order without extra communication between the peers~\cite{keidar23}.

In this research paper, we deal mainly with the following three axes of distributed systems: security, decentralization, and resilience.
In the first axis, security, we discuss the security of an open distributed system, its permission model, anonymity, confidentiality, and encryption.
Following the security axis, we explore decentralization as a way to circumvent centralized control and avoid single points of failure. Several design decisions in our system architecture influence the system's decentralization: a pure p2p system completely resides on the peer devices; while this approach highlights communication which is independent from any existing networking architecture, it introduces limitations on the range of communication, as it only allows for communication based on proximity due to the physical device constraints.
Finally, we discuss resilience, which in our context refers to network survivability under communication issues. These communication issues include Byzantine behavior, which is independent of the communication. These issues include reachability, bandwidth limitations, and coverage loss.

This paper proposes a system architecture for resilient unstructured p2p networks, based on an effective communication algorithm that leverages available computing power at the edge, to achieve grassroots consensus.
To evaluate the system architecture, we compare its central grassroots consensus communication algorithm against traditional distributed consensus algorithms: Paxos and Raft.
For this, in section~\ref{chap:background}, we discuss consensus algorithms for unstructured peer-to-peer networks, provide insights about their performance characteristics, and lay out their foundational equations.
Following the background chapter, section~\ref{chap:related} introduces related works and positions this research work contrastingly.
To formalize the setting, goals, and environment for this research paper, section~\ref{chap:problem} will introduce our basic assumptions, research question, and our approach to solve it.
Elaborating on the approach introduced in chapter~\ref{chap:problem}, section~\ref{chap:architecture} introduces the system architecture designed to support a permissionless, unstructured p2p social network, emphasizing the base algorithms, system components, and different communication layers.
The~\ref{chap:testbed} section introduces the testbed we used to evaluate the performance of the system architecture introduced in~\ref{chap:architecture}.
Next, chapter~\ref{chap:evaluation} evaluates the results of the experiment conducted on the testbed, which proves the possibility of our design.
Chapter~\ref{discussion} discusses the implications of the results introduced in the previous chapter, and finally, chapter~\ref{chap:conclusion} concludes the paper.